Napriek uvedeným komplikáciám bolo možné predpokladať, že v dlhých čítaniach bude možné identifikovať dostatočne dlhé opakujúce sa úseky, ktoré reprezentujú tú istú pôvodnú DNA sekvenciu.
Na základe tohto predpokladu bol identifikovaný klaster obsahujúci viac ako 37~000 sekvencií DNA.
Spočiatku sa zdalo, že ide o opakované čítanie jednej sekvencie, čo by mohlo viesť k získaniu kvalitného konsenzu.
Podrobnejšia analýza však ukázala, že spoločná zhoda medzi sekvenciami bola obmedzená iba na úvodný úsek s dĺžkou 43 nukleotidov.

Keďže 43 nukleotidov predstavuje príliš krátky úsek na spoľahlivú reprezentáciu celého DNA reťazca, no zároveň je príliš dlhý na náhodne sa vyskytujúci úsek, pôvodná hypotéza o prítomnosti dostatočne dlhého spoločného jadra nebola potvrdená.
Na základe tohto zistenia bola formulovaná nová hypotéza: spoločný úvodný úsek pravdepodobne zodpovedá DNA primeru prítomnému na začiatku sekvenovaných fragmentov.

%TODO: zformuj do akademického neosobného štýlu ale tak aby postup ale dalsi problem.
Nasledujúcim krokom preto bolo lokálne hľadanie primeru v jednotlivých čítaniach.
Primer sa nevyhľadával iba na začiatku reťazca, ale v celom rozsahu každého čítania, aby sa zachytili aj prípady jeho posunu alebo výskytu vo vnútri sekvencie.

